\section{Related Work}
\label{sec:related}

\textbf{Budgeted crawl and refresh ordering.} Pandey and Olston \cite{pandey2005usercentric} schedule the refreshing of stored web pages by assigning each page $p$ at time $t$ a priority $P(p, t)$ proportional to expected quality improvement and downloading the $B$ highest-priority pages each time unit. In follow-up work, Pandey and Olston \cite{pandey2008crawlorder} prioritize fetching new pages by search impact, defined as how often a page enters the top $K$ results for workload queries (for example, $K = 10$) and estimated from current top-$K$ workload scores.\looseness=-1

\textbf{Stale source summaries in federated search.} Kulkarni and Callan \cite{kulkarni2015selective} study selective search, which partitions a collection by similarity into topic shards and searches only a few shards per query. When collections change over time, routing summaries become stale \cite{ipeirotis2005changes,shokouhi2007updating}. Ipeirotis et al. \cite{ipeirotis2005changes} model content changes with survival analysis under limited network and compute resources to contact databases only when needed so summaries stay acceptably up to date. Shokouhi et al. \cite{shokouhi2007updating} evaluate forty-four dynamic collections over eight weeks, showing that out-of-date representations degrade accuracy and proposing three update policies: constant update (CU, expanding all representations equally each step, the uniform analogue of round-robin refresh), query-log popularity (QL), and collection size (SS).\looseness=-1

\textbf{Push synchronization and status updates under bandwidth limits.} Olston and Widom \cite{olston2002besteffort} study best-effort cache synchronization when bandwidth limits prevent exact synchronization, so stale copies are permitted. In their source-cooperation design, each source may wait for several updates to accumulate when bandwidth is limited, locally prioritizes its own modified objects, and refreshes modified objects whose priority exceeds its own estimate of a global refresh threshold $T$, which the cache regulates with occasional feedback messages. Maatouk et al. \cite{maatouk2020aoii} introduce the Age of Incorrect Information (AoII), a metric that addresses shortcomings of Age of Information and error penalty functions by extending fresh updates to fresh informative updates, so that the penalty grows the longer the monitor has incorrect information.\looseness=-1

\textbf{Partitioned vector index upkeep under drift.} Recent work maintains partitioned vector indexes as data and queries evolve. Xu et al. \cite{xu2023spfresh} show that updating one-third of the vectors in a static partition index degrades recall because fixed centroids miss gradual shift and raises tail latency because posting lists grow, and propose SPFresh to split partitions and reassign vectors in nearby partitions. Baranchuk et al. \cite{baranchuk2023dedrift} show content drift degrades accuracy without reconstruction and propose DeDrift-Lazy to update centroids from assigned vectors. Mohoney et al. \cite{mohoney2024adaivf} combine local re-clustering in Ada-IVF with a policy based on partition temperature, size deviation, and centroid drift, while Quake \cite{mohoney2025quake} splits or merges multi-level partitions using a query-latency cost model. Mageirakos et al. \cite{mageirakos2025cracking} propose CrackIVF, which starts small and progressively expands and adapts partitions as queries arrive. Jose \cite{jose2026repartition} studies when a partitioned approximate nearest-neighbor (ANN) index needs active re-partitioning under drift, showing that whether re-partitioning is needed is governed by the router's top-$b$ probe-budget slack (the number $b$ of cells each query probes, which plays the role of our $m$ across shards). Moderate centroid staleness can therefore be traded against a slightly larger probe budget. Single-index maintenance policies trade local maintenance work against extra partition probes within one index.\looseness=-1

\textbf{Distributed vector search routing and summary of differences.} Guo et al. \cite{guo2022manu} present Manu, where shards correspond to insertion and deletion channels, entities are hashed by primary key, and data enters growing segments that seal at a size or idle-time limit. Wang et al. \cite{semord2026} present Semord, where peers in a decentralized overlay exchange heartbeats with liveness state and compact synopses under a capped per-peer synopsis budget, gate epoch synthesis by hysteresis, and guard cached records with a time-to-live and epoch number. They also report refresh overhead in messages and bytes, and Semord peers contact extra peers via semantic expansion (the decentralized analogue of extra shard search) to compensate for routing error, whereas Semord's coordinator baseline uses fresh routing metadata. We instead study a coordinator-routed store whose summaries go stale, under a per-epoch byte budget $B$.\looseness=-1

Across these areas, we adapt top-$K$ search-impact prioritization \cite{pandey2005usercentric,pandey2008crawlorder}, source-cooperative local prioritization with coordinator cutoff feedback \cite{olston2002besteffort}, and AoII-motivated informative updates \cite{maatouk2020aoii} to $16$-center shard summaries. We compare that rule against the constant-update (round-robin), query-log popularity, and collection-size policies of Shokouhi et al. \cite{shokouhi2007updating} and a change-driven score that, like the schedules of Ipeirotis et al. \cite{ipeirotis2005changes}, favors summaries most likely to be out of date, but measures change at the shard instead of predicting it. Finally, we test an analogue of the probe-slack hypothesis of Jose \cite{jose2026repartition} and the extra-peer expansion of Semord \cite{semord2026} across distributed shards by comparing all five summary refresh rules against batch and causal online extra shard searches and coordinator vector push at equal budget $B$.\looseness=-1
